\documentclass[12pt]{article}

\usepackage[dvips,letterpaper,margin=0.75in,bottom=0.5in]{geometry}
\usepackage{cite}
\usepackage{slashed}
\usepackage{graphicx}
\usepackage{amsmath}
\usepackage{amssymb}
\usepackage{braket}
\usepackage{pythonhighlight}
\usepackage{mathtools}
\usepackage{tikz}
\begin{document}

\title{PHY 112L \\ Entropy of a 2-D Ideal-Gas}
\author{Michael Mulhearn}

\maketitle

\section{Introduction}

In this assignment, will will introduce the Gibbs Entropy as a
convenient way to measure changes in entropy for a numerical
simulation.  We will apply it to measure the change of entropy of an
ideal gas due to heating and free expansion in a molecular simluation,
and compare it to our expecation.  Last we will examine a gas
expanding into a region at a higher potential energy.

\section{Gibb's Entropy}

The entropy of a macrostate is defined as:
\begin{equation}
S = k \ln \Omega
\end{equation}
where $k$ is Boltzmann's constant and $\Omega$ is the multiplicity of
microstates corresponding to the macrostate.

To compute the Gibbs Entropy, we divide the total phase space into $M$
cells.  The macrostate is specfied by the number of particles in each
cell $n_i$.  To count these microstates, start by assuming the we have
only two cells, with the first containing $a$ particles.  The binomial
theorem gives us the multiplicity of states:
$$\Omega_{(2)} = \frac{N!}{a! (N-a)!}$$
If we divide the second cell into two cells containing $b$ and $c$ particles respectively, with:
$$N-a = b+c$$
then the binomial theorem implies that the multiplicity of states is:
$$\Omega_{(3)} = \frac{N!}{a! (N-a)!} \cdot \frac{(N-a)!}{b! c!} = \frac{N!}{a! b! c!}$$
If we continue in this way for $M$ cells, we arrive at the multinomial constant:
\begin{equation}
\Omega = \frac{N!}{\prod_{i=1}^M n_i!}
\end{equation}
Calculating
$$\ln \Omega = \ln N! - \sum_i \ln n_i!$$
and applying Stirling's approximation:
$$\ln \Omega = N \ln N - N - \sum_i (n_i \ln n_i - n_i)$$
Noting that:
$$N = \sum_i n_i$$
we have:
\begin{eqnarray*}
  \ln \Omega &=& -\sum_i n_i (\ln n_i - \ln N)\\
  &=& \- \sum_i n_i \ln \left( \frac{n_i}{N} \right)\\
\end{eqnarray*}
and so:
$$\frac{S}{kN} = \- \sum_i \frac{n_i}{N} \ln \left( \frac{n_i}{N} \right) $$
and defining:
$$p_i = \frac{n_i}{N}$$
we have at last:
\begin{equation}
\frac{S}{kN} = \- \sum_i p_i \ln ( p_i )
\end{equation}
which is known as the Gibbs Entropy (or Shanon Entropy for $k$=1).

As is always the case for classical entropy, the entropy is only
determined to within an overall constant, which depends on the cells
that we have choosen.  However, as long as we keep our choice of bins
constant, we can reliably calculate a change in entropy between
initial and final states using this formula.  To compute $p_i$, we
need only count the number of particles present in each cell and
divide by the total number of particles.

To determine the values of $p_i$ more reliably (with more statistics)
we can simulate many time steps, and count the number of visits of
particles to each cell.  Dividing the result by the total number of
visits (the number of particles times the number of time steps) we
determine $p_i$ in each cell with more statitics.  This only works
once the system has reached equilibrium (or else the Gibbs Entropy
will be changing).

\section{Entropy Changes for an Ideal Gas}

If no work is done on our gas, then:
$$dU = Q$$
and so
$$dS \equiv \frac{Q}{T} = \frac{NkdT}{T}$$
Thus the change in entropy for a gas that is heated at constant volume is given by:
$$\Delta S = \int_{T_I}^{T_f} \frac{NkdT}{T}$$
and so:
\begin{equation}
\Delta S = Nk \ln \left( \frac{T_f}{T_i}\right)
\end{equation}
Next, suppose that instead that we expand the gas from $V_1$ to $V_2$
at constant temperature.  The velocity profile remains unchanged, due
to the constant temperature, and so the phase-space only changes by
the volume, with:
$$\Omega \propto V^N$$
and so:
\begin{eqnarray*}
  \Delta S &=& k \, \ln(\Omega_2) - k \, ln(\Omega_1) \\
           &=& k \ln\left(\frac{\Omega_2}{\Omega_1}\right) \\
\end{eqnarray*}
and so:
\begin{equation}
  \Delta S = N \,k \, \ln \,\left( \frac{V_2}{V_1} \right)
\end{equation}

\section{Ideal Gas with a Step Potential}

Consider an ideal gas contained within $0 \leq y \leq 4 a$ subject to a potential energy:
\begin{equation}
  V(y) = \begin{cases}
    V_{\rm S} & y > 2 a \\
    0 & {\rm otherwise} \\
    \end{cases}
\end{equation}
The relative occupancy of the upper and lower regions (which are the same area) will be:
$$\frac{p_{\rm up}}{p_{\rm dn}} = \exp \left( - \frac{V_{\rm S}}{kT} \right) \equiv \alpha $$
and so:
\begin{equation}
p_{\rm up} = \frac{\alpha}{1 + \alpha}
\end{equation}
and
\begin{equation}
p_{\rm dn} = \frac{1}{1 + \alpha}
\end{equation}
If we consider this like a two state system going from the ground state ($S=0$) to a state at equilibrium with occupancies $p_{\rm up}$ and $p_{\rm dn}$, we expect the change in entropy to be:
$$\Delta S = p_{\rm up} \ln p_{\rm up} + p_{\rm dn} \ln p_{\rm dn}$$

\section{System of Units}

We will use the same system of units as the one we use for the first
assignment (as there is no longer any need to add a constraint for
gravity, which we will not consider in this assignment).  We assume
there is some relevant length scale $a$ which we set to $1$ and a
relevant temperature $T_0$ which we set to $1$ as well.

The other units that we defined in the previous homework continue to
apply, so let's recall them here.  Motivated by the Maxwell-Boltzmann
distribution, the velocity unit $v_0$ is defined by:
$$ v_0^2 \equiv \frac{k T_0}{m}$$
and the unit for time is:
$$ \tau \equiv a / v_0$$
Motivated by the ideal gas law, we can set the unit of pressure to:
$$O_0 \equiv \frac{NkT_0}{a^2}$$
and the unit of total energy:
$$U_0 = N k T_0$$
For this assignment, will measure potential energy (of a single particle) using units of
$$V_0 - k T_0$$
Remember that for analytic work, we keep these factors in place.  But when using an
analytic result within our computer programs, we set all of these quantities to 1:
$$m = a = g = T_0 = v_0 = \tau = O_0 = U_0 = V_0 = 1$$
As before, the temperature of an ideal gas is still calculated through
the average kinetic energy, which in our system of units becomes:
\begin{equation}
 \label{eqn:expv}
  \frac{T}{T_0} = \frac{\braket{v_x^2/v_0^2} + \braket{v_y^2/v_0^2}}{2}
\end{equation}
This is also the total kinetic energy of the gas, in units of total energy.

\section{Homework Problems}

{\bf Zero credit if you do not use our system of units consistently!}
When I ask you to calculate a quantity, e.g. temperature, I want the
answer in our units, e.g. $T_0$.  Do not first calculate in SI units
and then convert!\\[3pt]

\noindent
The starting point for this assigment is Problem 4 from the previous
homework.  You will want to initialize a numerical model for an ideal
gas tracking variables $v_x$, $v_y$, and $y$ for $N=10000$ molecules,
with height $H=4a$.  If you haven't already, define a function that
brings your gas into thermal contact with a reservoir of temperature
$T_R$, with an input parameter for the number of collisions.  Call it
what you want, but we will refer to it as {\tt collide\_reservoir}
below.  Use {\tt collide\_reservoir} to bring your gas into
equilibrium with a reservoir with $T_r = 1 T_0$.\\[3pt]

\noindent

{\bf Problem 1:} Write the function {\tt gibbs\_entropy} that
computes the Gibbs Entropy for a gas using the specified bins.
Validate your function on an initial distribution that is uniformly
distributed in $v_x$, $v_y$ and $h$.

{\bf Problem 2:} Calculate the the change in entropy when your gas is
heated from $T_I = 1$ to $T_I =2$ and compare to expectation.

{\bf Problem 3:} Simulate the free expansion of a gas from $h=2a$ to $h=4a$ at a contant temperature $T_I = 1$.  Compare your measured change in entorpy to expectation.

{\bf Problem 4:} Simulate the expansion of a gas from $h=2a$ to $h=4a$
at a contant temperature $T_I = 1$ when the upper region from $h=2a$
to $h=4a$ is at a higher potential energy $V_0$ compared to lower
region which you can take as $0$.  Show that the occupancy of the
upper region matches your expectation.

{\bf Problem 5:} Calculate the change in entropy from Problem 4 and
compare to expectation.

\end{document}
